% Source-only appendix candidate. Do not compile without the author's instruction.
% Based on the delivered CPII_Comparator_v2 workbook, inspected 2026-10-26.
\chapter{스프레드시트 시뮬레이터 사용 설명서}
\label{app:simulator-manual}

이 부록은 동봉한 \nolinkurl{CPII_Company_Comparator.xlsx}의 실제 화면을 기준으로 입력, 비교, 점검, 보관 절차를 설명한다. 대상 파일은 배포 폴더 \nolinkurl{outputs/client_delivery/CPII_Comparator_v2/}에 있으며, 파일 버전은 \nolinkurl{CPII_COMPARATOR_XLSX_V2.0.0}이다. 지수 산식과 해석은 부록~\ref{app:cpii_comparator_v2_method}를 참조하고, 여기서는 사용자가 어느 시트의 어떤 셀을 다루어야 하는지에 집중한다.

\begin{ReportCallout}
\textbf{처음 사용할 때는 동봉 예시로 계산 흐름부터 확인한다.} 비교기는 기업별 10개 변수를 받아 점수와 차이를 계산한다. 현재 등록된 비교기준 \nolinkurl{BM-DIGITAL-CONTENT-2025Q4}는 합성 예시용이다. 멜론과 스포티파이의 실제 자료를 적용하려면 서비스·시장·기간에 맞는 비교군, 변수 산식, 원천자료 해시를 등록하고 통합문서를 재생성해야 한다. 회사명만 바꾸어 기존 예시를 실제 기업의 결과로 제시해서는 안 된다.
\end{ReportCallout}

\section{빠른 시작: 열기에서 결과 확인까지}

\begin{enumerate}[leftmargin=1.8em,itemsep=0.45em]
\item 원본 파일을 보관하고, 날짜와 분석 목적을 붙인 작업 사본을 만든다. 예를 들어 \nolinkurl{CPII_Comparator_20261026_working.xlsx}로 저장한다. Excel의 계산 옵션은 자동으로 설정한다.
\item \nolinkurl{00_README}에서 버전과 입력 한도를 확인한다. 기본 예시는 \nolinkurl{01_INPUT_COMPANIES}의 7--10행에 들어 있다.
\item \nolinkurl{03_COMPARE_2}의 B4에서 \texttt{SYN-A}, B5에서 \texttt{SYN-B}를 선택한다. A14에 ``비교 가능''이 표시되는지 확인한다.
\item 같은 시트의 B8:D11에서 총점과 하위지수를 읽고, H17:H26에서 변수별 기여도를 확인한다. H28의 기여도 합계가 D8의 총점차와 일치하면 I28에 \texttt{PASS}가 표시된다.
\item \nolinkurl{06_SENSITIVITY}에서 사양을 바꿔도 차이의 방향이 유지되는지 살핀다. \nolinkurl{99_QA_AUDIT}에서 입력 오류와 합계 검증 상태를 확인한 뒤 결과를 저장한다.
\end{enumerate}

파란 입력 셀과 회사 선택 셀만 수정하는 것을 원칙으로 한다. 배포본에는 시트 보호가 설정되어 있지 않아 계산식도 덮어쓸 수 있다. 따라서 수식 영역을 직접 고치기보다 입력값을 수정하고 결과가 바뀌는지 확인한다. 원본 복사본은 수식 손상이나 입력 착오가 생겼을 때 되돌아갈 기준이 된다.

\clearpage
\section{시트 구성과 작업 범위}

\begin{longtable}{@{}L{4.0cm}L{3.1cm}L{7.6cm}@{}}
\caption{실제 배포 통합문서의 시트와 사용 위치}\label{tab:simulator-sheet-map}\\
\toprule
\rowcolor{ReportTableHead}
시트 & 작업 위치 & 용도 \\
\midrule
\endfirsthead
\toprule
\rowcolor{ReportTableHead}
시트 & 작업 위치 & 용도 \\
\midrule
\endhead
\bottomrule
\endfoot
\nolinkurl{00_README} & 읽기 전용으로 사용 & 빠른 시작, 버전, 점수의 의미, 입력 한도. \\
\nolinkurl{01_INPUT_COMPANIES} & A7:AC106 입력; AD열 점검 & 기업·서비스·시장·분기 정보와 활성 변수 10개. 최대 100개 관측. \\
\nolinkurl{02_INPUT_ECON_BASES} & A7:Z106 입력; AA열 점검 & 채널별 금액을 해석할 때 필요한 경제적 노출액과 가정. \\
\nolinkurl{03_COMPARE_2} & B4, B5 선택 & 두 기업의 총점, 하위지수, 변수별 점수차 기여도. \\
\nolinkurl{04_COMPARE_N} & A7:G106; J7:AM36 확인 & 비교집단별 순위와 최대 30개 관측의 짝비교 행렬. \\
\nolinkurl{05_ECON_INTERPRET} & 8--17행 확인 & 원값 차이와 노출액을 결합한 채널별 상당액·시나리오 금액. \\
\nolinkurl{06_SENSITIVITY} & A8:E11, B14:B18 확인 & 기본 사양과 세 가지 대체 사양의 점수차 비교. \\
\nolinkurl{07_INPUT_PROVENANCE} & A7:J1006 입력 & 관측별·변수별 출처, 산식, 품질, 이월 기록. \\
\nolinkurl{90_CALC_FEATURES} & 수정하지 않음 & 변수 변환과 정규화 계산. \\
\nolinkurl{91_CALC_SCORES} & 수정하지 않음 & 하위지수, 총점, 입력 검증 계산. \\
\nolinkurl{92_REF_VARIABLES} & 정의 확인 & 변수명, 운영 정의, 방향, 가중치, 경제적 해석 유형. \\
\nolinkurl{93_REF_ANCHORS} & 수정하지 않음 & 고정 정규화에 사용하는 기준값. \\
\nolinkurl{94_EXPECTED_OUTPUTS} & 수정하지 않음 & 원본 합성 예시의 기대 계산값. \\
\nolinkurl{99_QA_AUDIT} & A7:D18 확인 & 계산 일치, 입력 오류, 비교 가능 상태, 버전 점검. \\
\end{longtable}

계산 시트와 기준 시트도 화면에 표시된다. 이 시트들은 결과의 근거를 확인하기 위한 것으로, 일반 입력 화면처럼 사용하지 않는다. 106행 아래에 기업을 추가해도 수식과 검증 범위가 자동 확장되는 구조는 아니다.

\section{기업 정보와 10개 변수 입력}

\subsection{먼저 비교할 관측단위를 맞춘다}

\nolinkurl{01_INPUT_COMPANIES}의 한 행은 기업 전체가 아니라 하나의 \textbf{기업--서비스--시장--분기 관측}을 뜻한다. 7행부터 입력하며, A열의 관측 ID는 다른 행과 겹치지 않게 정한다. ID에는 영문 대문자, 숫자, 하이픈, 밑줄, 마침표를 1--64자 사용할 수 있다. B열 회사명에는 한국어를 사용할 수 있다.

D열 산업, E열 서비스 범위, F열 시장, G열 분기, H열 비교기준 ID가 모두 같은 관측끼리 비교한다. 영문 대소문자와 공백까지 일치해야 한다. 예를 들어 멜론의 한국 음악 구독 서비스와 Spotify 연결재무제표 전체를 같은 행 수준으로 놓지 말고, 같은 한국시장 서비스 범위와 기간에 맞춘 자료를 먼저 구성한다. 비교기업을 추가해도 고정 기준분포가 유지되므로 기존 기업의 점수는 바뀌지 않는다. 비교집단 안의 순위는 새 기업의 점수에 따라 달라질 수 있다.

\begin{longtable}{@{}L{2.0cm}L{4.2cm}L{8.5cm}@{}}
\caption{기업 입력 시트의 메타데이터}\label{tab:simulator-metadata}\\
\toprule
\rowcolor{ReportTableHead}
열 & 입력 내용 & 작성 방법 \\
\midrule
\endfirsthead
\toprule
\rowcolor{ReportTableHead}
열 & 입력 내용 & 작성 방법 \\
\midrule
\endhead
\bottomrule
\endfoot
A--C & 관측 ID, 회사명, 기업 구분 & C열은 \texttt{domestic}, \texttt{foreign}, \texttt{other} 중 선택한다. \\
D--H & 산업·서비스·시장·분기·비교기준 & 동일 비교집단의 다섯 필드를 통일한다. G열은 \texttt{2025Q4}와 같은 분기 형식으로 작성한다. \\
I--J & 통화, 가격연도 & 예: \texttt{KRW}, \texttt{2025}. 금액 입력과 같은 기준을 사용한다. \\
K--M & 품질, 합성 여부, 구축방식 & 품질은 A/B/C/\texttt{SYNTHETIC}, 합성 여부는 \texttt{yes}/\texttt{no}. 구축방식은 \texttt{direct}/\texttt{derived}/\texttt{proxy}/\texttt{mixed} 중 선택한다. 합성 예시는 품질과 합성 표식을 함께 유지한다. \\
N--P & 출처 ID, 출처일, 이월분기 & 출처일은 실제 날짜를 \texttt{YYYY-MM-DD} 텍스트로 입력한다. Excel 날짜 일련번호를 넣지 않는다. 이월분기는 0 이상의 정수이며 1을 넘으면 경고 대상이다. \\
Q--R & 노출조정 완료 여부, 노출계수 & 시장 공통값을 기업별 노출에 맞추어 조정한 뒤 Q열에 \texttt{yes}를 기록한다. R열은 0--1의 조정계수 기록이며, 입력변수에 자동으로 곱해지지 않는다. \\
S & 메모 & 서비스 범위, 자료 연결, 특이사항을 적는다. 선택 항목이다. \\
\end{longtable}

\subsection{변수는 이름과 단위를 확인해 입력한다}

T--AC열에는 10개 변수를 빠짐없이 입력한다. \textbf{1\%는 0.01, 3배는 3}으로 쓴다. 자료가 없는 칸을 0으로 채우면 ``부담이 0으로 관측되었다''는 다른 뜻이 되므로 빈칸으로 남기고 자료를 보완한다. 비교기는 10개 값이 모두 준비된 관측에 대해 점수와 순위를 제공한다.

\begin{longtable}{@{}L{0.8cm}L{4.8cm}L{9.1cm}@{}}
\caption{활성 변수의 실제 열 순서}\label{tab:simulator-features}\\
\toprule
\rowcolor{ReportTableHead}
열 & 변수 & 입력 시 확인할 내용 \\
\midrule
\endfirsthead
\toprule
\rowcolor{ReportTableHead}
열 & 변수 & 입력 시 확인할 내용 \\
\midrule
\endhead
\bottomrule
\endfoot
T & 유효세율 격차 & 기업 유효세율에서 동일 분야·분기 비교군의 기준값을 뺀 비율. 음수도 가능하다. \\
U & 인프라 비용·해외의존 결합비율 & 이 파일의 관계식은 V열 망비용률과 AB열 해외의존비율의 곱이다. \\
V & 네트워크 비용률 & 네트워크·클라우드·CDN 비용 등에 대한 대응 비용항목을 매출로 나눈 비율. \\
W & 세전이익 기반 회수 대리값 & 실행자료의 기본 정의는 세전이익/콘텐츠 투자액이다. 현금흐름 자체와 혼동하지 않는다. \\
X & 자금조달비용 스프레드 & 기업 조달비용률에서 같은 비교군 기준값을 뺀 비율. 노출액에 적용할 때 연율·분기율을 확인한다. \\
Y & 투자 회수율 & 기본 정의는 매출/콘텐츠 투자액이다. 배수로 입력한다. \\
Z & IP 자산화 결손 & 비교군의 무형자산/매출 기준값에서 기업의 해당 비율을 뺀 값. 세율 격차와 차감 방향이 다르다. \\
AA & 시장간 전이 대리지수 & 분야별로 구성된 0--1 범위의 대리값. 해당 분야의 산식과 자료 근거를 함께 기록한다. \\
AB & 해외의존비율 & 서비스·시장·분기에 대응하는 해외 트래픽 또는 외부 인프라 의존의 비율·대리값. \\
AC & 비용 대비 설비투자 압박률 & 네트워크·클라우드 비용/CAPEX의 배수. 실제 국외유출액을 직접 입력하는 칸은 아니다. \\
\end{longtable}

분야별 대리변수의 구체적 정의는 \nolinkurl{92_REF_VARIABLES}의 O열에서 확인한다. 투자 회수율과 세전이익 기반 회수 대리값의 부담 방향 전환은 계산기가 처리하므로 사용자가 미리 부호를 뒤집지 않는다. 입력이 끝나면 AD열의 검증 결과를 확인하고, 오류가 있는 행부터 바로잡는다.

\subsection{각 변수의 근거를 남긴다}

\nolinkurl{07_INPUT_PROVENANCE}에는 \textbf{관측 하나당 변수별 10행}을 작성한다. A열 관측 ID와 B열 변수명을 연결키로 쓰며, 입력 시트와 행 순서가 같을 필요는 없다. C--H열에는 출처 ID, 기준일, 품질, 구축방식, 이월분기, 비교기준 ID를 적고, I열에는 계산식과 가정을 남긴다. J열은 합성 여부다. 이 시트의 구축방식은 변수별로 \texttt{direct}, \texttt{derived}, \texttt{proxy} 중 하나를 선택한다.

출처 ID만 적어두면 원자료를 찾기 어려우므로, 같은 ID로 연결되는 자료 목록에 파일명·페이지·표 번호 또는 공개 URL을 보관한다. 가공값은 분자와 분모, 통화와 단위, 기준값을 적는다. 입력 검증이 통과했다는 사실은 자료의 형식이 맞는다는 뜻이며, 원자료와 산식의 대조는 별도로 수행한다.

\section{동봉 예시로 두 기업 비교를 재현한다}

가온콘텐츠(\texttt{SYN-A})와 누리미디어(\texttt{SYN-B})는 사용법 설명을 위해 만든 가상 기업이다. 입력 시트의 7행과 8행을 그대로 둔 상태에서 다음 순서로 결과를 확인한다.

\begin{enumerate}[leftmargin=1.8em,itemsep=0.45em]
\item \nolinkurl{03_COMPARE_2}의 B4에 \texttt{SYN-A}, B5에 \texttt{SYN-B}를 선택한다. A14의 상태가 ``비교 가능''이면 결과를 읽는다.
\item B8의 A 점수는 39.896, C8의 B 점수는 32.217, D8의 A--B는 7.678점이다. 이 비교에서는 A의 부담지수가 더 높다.
\item D9:D11의 규제·위험·시장 하위지수 차이는 각각 7.652, 8.296, 7.087점이다. 총점에서 세 하위지수의 비중은 각각 3분의 1이므로 총점차 기여는 각각 2.551, 2.765, 2.362점이다. 하위지수 차이 세 개를 그대로 더하지 않는다.
\item H17:H26을 확인하면 어느 변수에서 차이가 생겼는지 알 수 있다. 이 예시에서는 유효세율 격차가 약 1.674점, 세전이익 기반 회수 대리값이 약 1.596점 기여한다. H28의 합계는 D8과 일치한다.
\item \nolinkurl{04_COMPARE_N}에서 네 가상 기업의 같은 비교집단 순위를 확인한다. 다온스토리 42.767점, 가온콘텐츠 39.896점, 라온플레이 35.984점, 누리미디어 32.217점 순이다. \textbf{1위는 부담지수가 가장 높은 관측}을 뜻한다.
\end{enumerate}

짝비교 행렬에서는 행 기업의 점수에서 열 기업의 점수를 뺀다. 기업 순서를 바꾸면 차이의 부호도 바뀐다. 다른 산업·서비스·시장·분기·비교기준에 속한 조합은 \texttt{N/A}로 표시되므로 전체 기업을 하나의 순위로 해석하지 않는다. 동점은 평균순위로 표시한다.

\subsection{입력값을 바꾸어 보는 연습}

원본과 구분한 작업 사본에서 가온콘텐츠의 V7 망비용률을 0.180에서 0.160으로 바꾸고, U7 결합비율도 \(0.160\times0.740=0.1184\)로 함께 바꾼다. AB7의 해외의존비율은 0.740을 유지한다. 메모와 변수별 출처 기록에 ``사용법 연습: 망비용률 가정 변경''이라고 남긴 뒤, AD7과 비교 화면을 다시 확인한다. 두 관련 변수의 정규화 위치와 기여도가 어떻게 달라지는지 원본 사본과 대조한다.

이 연습은 입력값에 대한 반응을 확인하는 절차다. 원값이 변해도 같은 기준 순위 구간 안에 있으면 점수가 그대로일 수 있다. 또한 원본 예시의 10개 값이 달라졌으므로 원본 기대점수 대조 대상에서 제외된다. 정상 입력에 \texttt{READY}가 표시되고 원본 대조 행수가 줄어드는 것은 이 변경에 따른 정상적인 동작이다.

\section{경제적 노출액을 입력하고 금액을 읽는 방법}

점수 비교만 할 때는 \nolinkurl{02_INPUT_ECON_BASES}를 비워두어도 된다. 채널별 금액 해석이 필요하면 같은 관측 ID로 경제적 노출액을 입력한다. 매출과 부채 등은 통화의 기본 단위로 작성한다. \texttt{KRW}라면 원 단위이므로 1조원은 \texttt{1000000000000}이다.

\begin{longtable}{@{}L{2.3cm}L{4.0cm}L{8.4cm}@{}}
\caption{경제적 노출액 시트의 입력 범위}\label{tab:simulator-economic-input}\\
\toprule
\rowcolor{ReportTableHead}
열 & 항목 & 입력 기준 \\
\midrule
\endfirsthead
\toprule
\rowcolor{ReportTableHead}
열 & 항목 & 입력 기준 \\
\midrule
\endhead
\bottomrule
\endfoot
A--E & 관측 ID·통화·가격연도·기간·연환산계수 & 통화와 가격연도는 기업 입력과 맞춘다. D열 기간은 해당 분기 또는 그 연도를 사용한다. E열은 실제 기간에 맞는 0 초과 1 이하의 값으로 입력하며, 연간은 1, 1개 분기는 필요에 따라 0.25를 쓴다. \\
F--H & 매출액 & 하한(Low), 기준(Base), 상한(High). 네트워크 비용과 IP 자산화 해석 등에 사용한다. \\
I--K & 양의 세전이익 & 세율 격차의 조건부 세부담 상당액에 사용한다. \\
L--N & 평균 이자부부채 & 조달비용 차이와 기간을 결합하는 기준이다. \\
O--Q & 콘텐츠 투자액 & 회수 관련 변수의 성과 차이를 읽는 기준이다. \\
R--T & 전이대상 매출 & 시장전이 시나리오의 대상 범위를 정한다. \\
U--W & 시나리오 전환율 & 대상 매출과 함께 입력하는 전환율 가정이다. \\
X--Z & 출처·품질·메모 & 노출액 근거, 기간 처리, 하한·상한의 설정 이유를 남긴다. \\
\end{longtable}

각 하한·기준·상한 묶음은 모두 입력하거나 모두 비운다. 값은 0 이상이며 하한\(\leq\)기준\(\leq\)상한이어야 한다. 전이대상 매출과 전환율은 두 묶음을 함께 입력한다. 연환산계수는 입력 기간을 자동 판별해 채우는 값이 아니다. 특히 조달비용률이 이미 분기 기준으로 환산되어 있다면 같은 기간계수를 다시 적용하지 않도록 원자료의 단위와 계산식을 먼저 맞춘다. 완료 후 AA열의 검증 상태를 확인한다.

\nolinkurl{05_ECON_INTERPRET}는 두 기업 선택을 \nolinkurl{03_COMPARE_2}에서 가져온다. I열의 상태가 산출 가능함을 나타내는지 먼저 확인하고, M--O열의 공통규모 금액과 P--R열의 실제규모 금액을 구분해 읽는다. 공통규모는 두 기업의 평균 노출액을 적용하여 원값 차이를 비교하고, 실제규모는 각 기업의 노출액을 적용하므로 규모 차이도 함께 반영한다. 각각의 가운데 열 N과 Q가 기준(Base)값이다.

\begin{ReportCallout}
\textbf{망비용 예시의 금액은 다음과 같이 읽는다.} 동봉 자료에서 A와 B의 망비용률은 0.180과 0.105, 기준 매출액은 각각 1조원과 1.4조원이다. 공통 매출규모 1.2조원을 적용한 N10은 \((0.180-0.105)\times1.2\text{조원}=900\text{억원}\)이다. 실제규모를 적용한 Q10은 \(0.180\times1\text{조원}-0.105\times1.4\text{조원}=330\text{억원}\)이다. 두 값은 규모를 처리하는 방식이 다른 같은 채널의 예시이며, 7.678점 전체를 환산한 금액이 아니다.
\end{ReportCallout}

인프라 결합값, 해외의존비율, 비용 대비 설비투자 압박률은 별도 금액을 산출하지 않는다. IP 자산화와 시장전이 금액은 입력 가정에 따른 시나리오로 읽는다. 채널 간 같은 비용이나 투자액을 공유할 수 있어 금액 총합은 제공하지 않는다. 결과를 옮길 때도 채널, 노출액 기준, 기간, 하한·기준·상한을 함께 표시한다.

\section{민감도와 시나리오의 사용 범위}

\nolinkurl{06_SENSITIVITY}는 사용자가 정책계수를 움직이는 화면이 아니라 \textbf{미리 정한 네 가지 계산 사양의 비교표}다. 8행은 기본 10변수 사양, 9행은 규제·위험·시장 비중 40/40/20, 10행은 시장간 전이·해외의존의 저지원 변수 제외, 11행은 인프라 결합값·비용 대비 설비투자 압박률·세전이익 기반 회수 대리값을 제외한 중복채널 축소 사양이다.

기본 예시 A--B의 점수차는 네 사양에서 각각 7.678, 7.796, 8.788, 6.399점이다. B14:B16에는 최솟값 6.399, 최댓값 8.788, ``양(+) 안정''이 표시된다. 이는 이 네 사양에서 차이의 방향이 같다는 뜻이다. 하한·상한은 신뢰구간이 아니며, B18의 ``미평가''는 입력변수별 범위 민감도를 이 Excel 화면에서 계산하지 않았다는 표시다.

입력변수의 하한·기준·상한을 바꾸는 별도 분석은 동봉 \nolinkurl{CPII_Feature_Range_Template.csv}와 Python 비교 엔진의 \nolinkurl{--feature-ranges} 옵션을 사용한다. 경제적 노출액의 Low/Base/High와 입력변수 자체의 범위는 서로 다른 가정이다. 또한 본문에서 제안한 \nolinkurl{SCN_BUNDLE}, \nolinkurl{SCN_NETWORK}, \nolinkurl{SCN_IPFIN}, \nolinkurl{SCN_BARGAIN} 탭과 정책계수 조절표는 이 배포 통합문서에 구현되어 있지 않다. 본문의 정책 시나리오 결과와 기업 비교기의 민감도 결과를 별도 표로 관리한다.

\clearpage
\section{빈칸·경고·오류가 나타날 때}

\begin{longtable}{@{}L{4.2cm}L{5.0cm}L{5.5cm}@{}}
\caption{주요 표시와 해결 순서}\label{tab:simulator-troubleshooting}\\
\toprule
\rowcolor{ReportTableHead}
표시 또는 현상 & 뜻 & 확인할 곳 \\
\midrule
\endfirsthead
\toprule
\rowcolor{ReportTableHead}
표시 또는 현상 & 뜻 & 확인할 곳 \\
\midrule
\endhead
\bottomrule
\endfoot
두 기업 비교 결과가 빈칸 & 선택 누락, 동일 ID 선택, 중복 ID, 입력 오류 또는 비교집단 불일치일 수 있다. 빈칸은 0점이 아니다. & 먼저 \nolinkurl{03_COMPARE_2}!A14, 다음으로 입력 시트 AD열을 확인한다. \\
짝비교 행렬의 \texttt{N/A} & 다섯 비교집단 필드가 일치하지 않는다. & 산업·서비스·시장·분기·비교기준의 값과 공백·대소문자를 대조한다. \\
인프라 관계식 경고 & 결합비율이 망비용률\(\times\)해외의존비율과 허용오차 이상 다르다. & U열, V열, AB열의 정의·단위·계산을 함께 확인한다. \\
\texttt{READY} & 입력이 계산 가능한 상태다. 수정된 예시나 새 관측은 원본 기대점수 대조 대상이 아닐 수 있다. & 원본 예시의 \texttt{PASS}와 구분하되, \texttt{READY} 자체를 오류로 보지 않는다. \\
\texttt{CHECK} 또는 원본 대조 불일치 & 원래 예시 입력과 대응하는 기대값의 계산이 맞지 않을 수 있다. & 수식을 덮어썼는지 확인하고, 변경하지 않은 원본으로 다시 대조한다. \\
\texttt{MISSING EXPOSURE} / \texttt{MISSING ASSUMPTION} & 해당 금액 계산에 필요한 노출액 또는 가정이 없다. & 경제적 노출액의 관련 3개 값과 전환율을 확인한다. 점수 계산과는 구분한다. \\
\texttt{CURRENCY/YEAR MISMATCH} / \texttt{TIME BASIS MISMATCH} & 통화·가격연도 또는 기간이 맞지 않는다. & 기업 입력과 경제적 노출액 입력, 두 비교기업의 기준을 맞춘다. \\
\texttt{NON-MONETIZABLE} & 이 변수에는 독립적인 금액환산식이 제공되지 않는다. & 빈 금액을 추정하여 채우지 말고 점수와 원값을 보고한다. \\
원값을 바꾸어도 점수가 같음 & 고정 기준분포의 같은 순위 구간에 있거나 상·하단 경계가 적용되었을 수 있다. & 자동 계산 상태와 입력 검증을 먼저 확인한 뒤 정규화 위치를 살핀다. \\
\texttt{\#REF!}, \texttt{\#VALUE!} 등 & 수식 참조, 데이터 형식 또는 계산 환경의 문제다. & 행·열 삭제, 문자형 숫자, 수식 덮어쓰기를 점검한다. 원본 사본으로 복원하는 편이 안전하다. \\
\end{longtable}

\nolinkurl{99_QA_AUDIT}에서는 D9의 기여도 합계 점검, D16의 기업 입력 점검, D17의 경제적 노출액 점검, D18의 선택기업 비교 상태를 확인한다. D7은 원본 예시와 Python 기준값의 수치 일치를 확인하는 항목이다. 예시를 수정하면 B8의 원본 대조 가능 행수가 줄 수 있다. 따라서 원본 대조 행수 감소와 입력 차단 오류를 구분한다.

\section{저장·복원·재현 절차}

결과를 공유할 때는 계산표 한 장과 함께 \textbf{그 결과를 만든 입력 사본}을 남긴다. 파일명에는 분석일과 시나리오명을 넣고, 다음 정보를 같은 폴더에 보관한다.

\begin{itemize}[leftmargin=1.8em,itemsep=0.35em]
\item 사용한 통합문서 버전, 방법 버전, 기준분포 버전과 두 관측 ID.
\item 기업 입력, 변수별 출처, 경제적 노출액, 변경한 가정과 변경 사유.
\item 비교 가능 상태, 총점차와 변수별 기여도, 민감도 결과, 주요 검증 상태.
\item 원자료 사본 또는 접근 경로와 자료 기준일. 통화·단위·기간의 변환 근거.
\end{itemize}

연습을 끝낸 뒤 초기 상태로 돌아가려면 수정하지 않은 배포 원본에서 새 사본을 만든다. 수식까지 포함된 시트 전체를 지우거나, 다른 파일에서 행을 통째로 덮어쓰는 방식은 피한다. 입력값을 새로 붙여넣을 때는 해당 입력 범위에 값만 붙여넣고 드롭다운과 서식이 유지되는지 확인한다.

배포 폴더에는 기업 입력, 변수별 출처, 경제적 노출액, 입력범위를 위한 CSV 템플릿과 재현 코드가 함께 있다. 100개 관측을 넘는 분석이나 화면 한도 30개를 넘는 전체 짝비교 행렬은 Python 엔진의 CSV 출력으로 관리한다. 배포본의 무결성은 해당 폴더에서 다음 명령으로 확인할 수 있다.

\begin{quote}\small\ttfamily
python3 verify\_release.py --verify .
\end{quote}

이 검사는 배포 파일과 기록된 해시의 일치를 확인한다. 입력을 바꾼 작업 사본은 원본과 파일 내용이 달라지므로 원본 배포 폴더와 분리해 보관한다. Excel에서 수정한 내용이 동봉 CSV에 자동 저장되지는 않는다. Python으로 같은 분석을 재현하려면 관측 ID와 변수명을 기준으로 변경 입력을 CSV에도 반영하고 검증해야 한다.

\begin{ReportCallout}
\textbf{제출용 결과표에는 숫자의 읽는 법을 함께 남긴다.} 표 제목에 서비스·시장·분기와 비교기준을 쓰고, A--B의 차감 순서와 ``값이 클수록 부담지수가 높음''을 명시한다. 금액은 해당 채널의 노출액·기간·가정을 덧붙인다. 이 네 가지가 갖추어지면 같은 계산을 재현하고, 새로운 자료가 들어왔을 때 무엇이 달라졌는지 확인할 수 있다.
\end{ReportCallout}
